\documentclass[conference]{IEEEtran}
\usepackage[utf8]{inputenc}
\usepackage[T1]{fontenc}
\usepackage{url}
\usepackage{booktabs}
\title{Geospatial Context for Residential Property Valuation: An H3-Enriched Gradient-Boosting Framework}
\author{\IEEEauthorblockN{Alvaro Careli}\IEEEauthorblockA{Independent Researcher\\Brazil}}
\begin{document}
\maketitle
\begin{abstract}
Property values reflect both property attributes and spatial context. This paper presents a reproducible valuation framework that enriches listing records with H3-indexed urban features: accessibility to services, POI counts, Sentinel-2 NDVI/NDBI statistics, zoning, and explainable opportunity scores. The implementation uses LightGBM with encoded categorical variables and a cross-validation plus hold-out evaluation workflow. We provide a rigorous experimental protocol contrasting intrinsic-only and geospatially enriched models, including spatial and temporal validation to prevent information leakage. No unverified accuracy values are reported: the available implementation establishes a testable pipeline, and the manuscript specifies the evidence required for a publishable empirical claim.
\end{abstract}
\begin{IEEEkeywords}automated valuation model, real estate, geospatial machine learning, LightGBM, H3, explainable AI\end{IEEEkeywords}

\section{Introduction}
Automated valuation models traditionally combine area, rooms, bathrooms, parking, property type, and location. Hedonic approaches formalize the relationship between price and property characteristics \cite{anand2016hedonic}, but coarse location variables can omit the urban context that buyers and businesses experience: nearby services, accessibility, land-use patterns, environmental conditions, and legal restrictions.

This paper describes an enrichment-first valuation architecture. Each listing is joined to an H3 cell, then augmented with spatial features produced by a territorial intelligence pipeline. A LightGBM regressor is trained on the enriched table. LightGBM is an efficient gradient-boosting decision-tree approach \cite{ke2017lightgbm}; nevertheless, model sophistication is not a substitute for a leakage-safe evaluation. Our contribution is a replicable experimental design that tests whether geospatial context adds out-of-sample value beyond intrinsic housing attributes.

\section{Architecture}
The input table contains price and available intrinsic attributes: area, bedrooms, bathrooms, parking spaces, property type, neighbourhood, and coordinates. If an H3 identifier is absent, latitude and longitude are converted to an H3 resolution-8 cell. The enrichment join supplies NDVI/NDBI means and slopes, POI counts, distances to supermarkets/pharmacies/schools, zone, and residential/commercial opportunity scores.

The model supports numeric variables and categorical property type, neighbourhood, and zone. Categorical fields use an ordinal encoder with a reserved unknown representation. Hyperparameters are selected with five-fold cross-validation, then a final estimator is fit on the training partition. Mean absolute error (MAE), $R^2$, cross-validated MAE mean, and cross-validated MAE standard deviation are recorded. Model and preprocessing metadata are stored together to ensure inference uses the training feature schema.

\section{Hypotheses and Experimental Design}
The main hypothesis is $H_1$: adding spatial features reduces prediction error relative to an intrinsic-only model on an untouched evaluation set. The study should compare:
\begin{itemize}
\item M0: intrinsic attributes only;
\item M1: M0 plus latitude/longitude and neighbourhood;
\item M2: M1 plus POI and accessibility features;
\item M3: M2 plus NDVI/NDBI, zoning, and opportunity scores.
\end{itemize}
All models must use identical records and a fixed preprocessing definition. Report MAE, median absolute error, RMSE, $R^2$, and the relative MAE reduction from M0. Use paired bootstrap confidence intervals for differences in absolute error, not just point estimates.

Random splits alone can overestimate performance when nearby listings are similar. The primary validation should group folds by neighbourhood or H3 parent cell. If listing dates are available, a second evaluation must train on earlier listings and test on later listings. Features must be computed using only information available on or before each listing date; otherwise POI changes, later satellite data, and aggregate scores can leak future information.

\section{Interpretation and Feature Governance}
The enrichment design produces useful features but does not imply causality. A distance to a school, for example, may correlate with price through unobserved neighbourhood characteristics. Model explanations such as SHAP values can communicate local feature contributions \cite{lundberg2017shap}, but must be framed as associations.

Every feature should carry source, spatial resolution, and observation date. OSM data should be audited for geographic coverage \cite{haklay2010osm}; Sentinel-2's spectral and spatial properties are documented by ESA \cite{esaS2}. Zoning must be treated cautiously: it can be informative, but it may encode planning classifications correlated with price. A valuation system must not use protected personal attributes or proxies intended to discriminate against residents or neighbourhoods.

\section{Current Implementation Status}
The software implementation contains the full enrichment interface, configurable feature lists, LightGBM training, Optuna hyperparameter tuning, five-fold cross-validation, a hold-out split, and optional MLflow logging. A small listing file is present for technical testing, but no held-out metrics are presented here because the dataset's provenance, temporal coverage, and target quality have not yet been established for scientific evaluation.

The paper should be updated only after assembling a documented dataset of completed transactions or consistently verified listings. Minimum reporting includes the number of records, city and period, missing-value policy, geographic distribution, price currency, all model hyperparameters, and results for M0--M3. Insert measured results here before submission: sample size = \texttt{TBD}; M0 MAE = \texttt{TBD}; M3 MAE = \texttt{TBD}; spatial-fold MAE = \texttt{TBD}.

\section{Conclusion}
We introduced an H3-enriched property valuation workflow that connects urban intelligence with gradient-boosting regression. Its technical contribution is a reproducible pipeline and a strict evaluation plan for testing incremental value from geospatial context. The next research milestone is not additional model complexity; it is a time-stamped, representative market dataset and spatially honest benchmark results.

\bibliographystyle{IEEEtran}
\bibliography{references}
\end{document}
